\documentclass[12pt,a4paper]{article}

\usepackage{cmap}
\usepackage[T2A]{fontenc}
\usepackage[utf8]{inputenc}
\usepackage[english,russian]{babel}
\usepackage[a4paper,left=2.5cm,right=2cm,top=2cm,bottom=2cm]{geometry}
\usepackage{booktabs}
\usepackage{array}
\usepackage{enumitem}
\usepackage[table]{xcolor}
\usepackage{hyperref}

\definecolor{hamster}{RGB}{240,148,10}
\definecolor{inkbrown}{RGB}{74,52,19}

\hypersetup{
  colorlinks=true,
  linkcolor=hamster,
  urlcolor=hamster,
  pdftitle={Тапаем хомяка — правила игры},
  pdfauthor={Тапаем хомяка}
}

\setlength{\parindent}{0pt}
\setlength{\parskip}{6pt}

\newcommand{\term}[1]{«\textbf{#1}»}

\title{\textbf{\textcolor{inkbrown}{Тапаем хомяка}}\\[4pt]
\large правила игры-соревнования}
\date{}

\begin{document}
\maketitle
\vspace{-2.5em}

\section*{Легенда}

Идёт 2052 год. Мировая экономика окончательно переехала в телеграм:
единственная твёрдая валюта человечества --- \emph{хомякоин}, который
добывается тапаньем по хомяку. Полтора миллиарда человек тапают
не поднимая головы, большие пальцы внесены в список стратегических
ресурсов, а аирдроп --- \emph{вот-вот, уже завтра, ну максимум через
квартал} --- обещают четырнадцатый год подряд.

Вы --- совет директоров хомячьей криптобиржи. Ваш хомяк круглосуточно
натапывает монеты, но у конкурентов хомяки не хуже, а у некоторых,
по слухам, хомяк с двумя большими пальцами. Единственный путь
к господству --- \emph{карточки апгрейдов}: их разрабатывают
учёные-хомяководы, они же --- задачи с сайта informatics. Внедрённый
апгрейд приводит в восторг криптоинвесторов (\emph{«Хомяк на солнечных
батареях?! Берём!»}), и они засыпают биржу монетами и ускоряют тапы.

Правда, пока лаборатория возится с апгрейдом, она нагло отвлекает
хомяка от тапанья --- добыча замедляется. А за отдельную плату знакомый
кит готов слить \emph{инсайд} с чужой биржи (тест к задаче). Кит берёт
предоплату, работает без гарантий, а на вопросы отвечает
«скоро листинг, брат».

К дедлайну, объявленному Всемирной Комиссией по Хомяководству, побеждает
биржа с самым большим счётом монет. Остальных ждёт слияние, поглощение
и переработка в скам-проект с обещанием аирдропа.

\section*{Как всё устроено}

\begin{enumerate}[leftmargin=1.6em,itemsep=2pt]

\item \textbf{Стартовый капитал.} Каждая биржа начинает с одинаковым
счётом \term{Монет} и одинаковой скоростью \term{Тапов в секунду}.
Хомяк не останавливается ни на секунду: каждую секунду игрового времени
счёт пополняется на текущую скорость тапов. Даже если вся команда ушла
смотреть график к луне.

\item \textbf{Каталог апгрейдов.} Всем биржам доступен один и тот же
набор задач: несколько уровней сложности, в каждом уровне --- несколько
задач. Но порядок задач внутри уровня у каждой биржи свой, секретный
и перемешанный: подглядывать в экран конкурента бессмысленно ---
их карточка № 5 и ваша карточка № 5 --- почти наверняка разные задачи.
Инсайд через плечо не работает, кит подтверждает.

\item \textbf{Покупка апгрейда (задачи).} Условие задачи засекречено,
пока апгрейд не куплен: платите \term{Цену задачи} со счёта ---
получаете ссылку на условие. Купить задачу команда может \emph{сама},
кнопкой на своей странице. Возвратов нет: потраченные монеты
не возвращаются \emph{никогда} --- это же крипта, привыкайте.

\item \textbf{Плата за прогресс.} С момента покупки и до сдачи апгрейда
лаборатория отвлекает хомяка: \term{Тапы в секунду} снижаются
на \term{Нагрузку от задачи}. Хотите купить всё и сразу --- считайте:
если монет на счету меньше цены, а хомяк и так еле тапает,
биржа сделку не одобрит.

\item \textbf{Инсайд от кита (тесты).} Любое количество,
по \term{Цене подсказки} за штуку, приобретается через ведущего ---
кит общается только с директором лично. Монеты за инсайд
не возвращаются, качество инсайда обсуждению не подлежит
(«не финансовый совет», --- добавляет кит).

\item \textbf{Внедрение апгрейда (решение задачи).} Решение сдаётся
на informatics. Как только вердикт --- \emph{OK} или \emph{Зачтено},
биржа автоматически получает всё и сразу:
\begin{itemize}[itemsep=0pt]
	\item счёт пополняется на \term{Бонус запасы} --- инвесторы ликуют;
	\item хомяк возвращается к тапанью
	(\term{Нагрузка} снова прибавляется к скорости);
	\item и сверху --- \term{Бонус скорость}: пальцы хомяка разгоняются
	навсегда.
\end{itemize}
Моментом внедрения считается \emph{время отправки посылки}, а не время
её проверки. Если тестирующая система задумалась --- это её просадка,
а не ваша.

\item \textbf{Посылка на флажке.} Решение, отправленное до дедлайна,
но проверенное после, --- засчитывается. Отправленное после дедлайна ---
уходит в архив Комиссии по Хомяководству с пометкой
\emph{«надо было раньше»}.

\item \textbf{Без самодеятельности.} Сдавать некупленную задачу
бесполезно: апгрейд, установленный без оплаты, автоматика не внедряет,
а отправляет ведущему на разбирательство. Ведущий, конечно, может
и зачесть --- но сперва очень выразительно на вас посмотрит.

\item \textbf{Победа.} В момент дедлайна побеждает биржа с наибольшим
счётом \term{Монет}. Скорость тапов, амбиции и обещания аирдропа
в зачёт не идут --- только монеты на счету.

\item \textbf{Ведущий --- окончательный авторитет.} Он видит все сделки
в журнале, может продлить игру, исправить ошибочную запись и вручную
зачесть решение, если informatics прилёг отдохнуть. Спорить с ведущим
можно, но статистика споров не на вашей стороне: 0 побед из 0 возможных
у предыдущих составов.

\end{enumerate}

\section*{Параметры уровней}

Цены и бонусы зависят от уровня сложности задачи и объявляются
на табло игры (таблица \emph{«Параметры игры»}). Чем сложнее апгрейд,
тем дороже инсайд и тем щедрее инвесторы.

\begin{center}
\begin{tabular}{lccc}
\toprule
 & \textbf{Уровень 1} & \textbf{Уровень 2} & \textbf{Уровень 3} \\
\midrule
Цена задачи        & 12\,000 & 12\,000 & 12\,000 \\
Цена подсказки     & 3\,000  & 7\,000  & 10\,000 \\
Нагрузка от задачи & 2       & 2       & 2       \\
Бонус запасы       & 12\,000 & 25\,000 & 50\,000 \\
Бонус скорость     & 4       & 7       & 11      \\
\bottomrule
\end{tabular}

\small\emph{(значения по умолчанию; точные параметры вашей игры --- на табло)}
\end{center}

\section*{Шпаргалка по цветам табло}

\begin{center}
\begin{tabular}{ll}
\toprule
\colorbox{gray!30}{\strut ~серая~}   & задача не куплена: карточка в сейфе, условие скрыто \\
\colorbox{yellow!60}{\strut ~жёлтая~} & куплена: лаборатория работает, хомяк страдает \\
\colorbox{green!40}{\strut ~зелёная~} & решена: инвесторы аплодируют стоя \\
\bottomrule
\end{tabular}
\end{center}

\vspace{1em}
\begin{center}
\emph{Удачного тапанья! И помните: возвратов нет, аирдроп завтра.}
\end{center}

\end{document}
